\chapter{The accounting language}
\label{ch:language}

\section{Scope of the first language}

The historical AccDSL Core 0 proposal describes entities and books, accounts, observations, shared
unknowns, constraints, alternative interpretations, contract instances,
accounting rules and reporting mappings. It has pure expressions and a
restricted posting effect. Economic forecasting algorithms can live in typed
plugins; unrestricted numerical code does not need to be reimplemented in
the DSL. The historical grammar in Appendix~\ref{app:contracts} fixes the preserved
teaching examples. The proposed portable records and first company profile
below define the new authoring boundary; the first matched text frontend
supports the declared debt fragment.

A subsequent documentary subset authors the funded-year extension of that
company profile: identified claim allocations, explicit credit terms and a
dated event sequence. Text and typed records lower to the same body, which
drives the conditional reference calculation in
Section~\ref{sec:funded-year-authoring}. The imported annual projection still
supplies the other operating and accounting rules. This extends the matched
example without establishing a complete whole-company text grammar.

The next subset makes the operating inputs explicit as well.
Section~\ref{sec:company-input-authoring} starts from the historical opening,
then derives purchases, sales, claims, expenses and closing statements from
authored quantities, prices, asset cohorts and contract terms. It imports no
forecast journals. This makes a price or service-date change traceable through
the accounts, while retaining the example's stated method restrictions.

Identified operations can then repeat within a year and carry their unpaid
claims and asset cohorts into the next year. The conditional path in
Section~\ref{sec:company-operation-paths} follows Northstar through three
annual periods, including revised assumptions and an interrupted year.
Its typed records specify this composition. The bounded text frontend in
Section~\ref{sec:company-path-authoring} preserves those records exactly,
including named settlement allocations and incomplete paths.
Section~\ref{sec:company-family-language} adds named unweighted scenarios and
pure typed expressions for the declared collection policy.
Section~\ref{sec:company-module-language} preserves their portable identities
and specifies the additional state needed by a shared-book return method.
Section~\ref{sec:company-return-reference} gives its matched authoring
and finite shared-book execution reference. General functions,
endogenous decisions, source admission and broader accounting profiles remain
language tasks.

The historical Core 0 compiler has six ordered stages: parsing with source locations; name and
module resolution; type and temporal-context checking; lowering to the IR;
constraint and rule-applicability analysis; and execution-plan construction.
The compiler never silently supplies an economic assumption to make a program
compile. A program can be valid as a historical representation while requiring
additional assumptions before a particular simulation can run.
The proposed portable compilation sequence is specified separately in
Section~\ref{sec:portable-authoring}.

The language's meaning is defined by its IR and evaluator, rather than by the
surface spelling alone. Formatting changes cannot change a model's semantic
hash. A changed source span or reviewer decision changes evidence provenance
even if the numerical model remains the same. The system therefore maintains
separate evidence, semantic-model and execution identifiers.

For the preserved Core 0 examples, the complete grammar is supplied as
\code{spec/core-0.lark}; the original design directory retains
\code{spec/core-ir.schema.json} and \code{spec/core-contract.md}. They define reference resolution, expression arity,
type/context checks, operations, failures and identity. Structural schema
validation alone does not establish those semantic properties. ASCII
identifiers and UTF-8 strings are used in this revision; currency suffixes,
dates and decimal literals have explicit lexical rules. The grammar is checked
by parsing the complete demonstration file, not by claiming that its parser
also implements a compiler or solver.

\section{Types that protect meaning}

The following table uses the preserved Core 0 type vocabulary. Its distinctions
motivate the more precise portable types introduced below.

\begin{longtable}{L{40mm}L{96mm}}
\toprule Type & Meaning and checks\\\midrule\endhead
\code{Money<C>} & Exact amount in currency C. Addition requires identical
currency; conversion requires a dated rate and policy.\\
\code{Quantity<U>} & Units of shares, goods or other commodities. Quantities
cannot be added to money or used as a functional-currency posting.\\
\code{Rate<B>} and \code{YearFraction} & Rate with basis B, plus a separately
computed accrual fraction. Simple and compound accrual functions are distinct.\\
\code{Instant}, \code{Period}, \code{Calendar} & Instant observations and
flows over intervals; interval operations preserve calendar conventions.\\
\code{Account<K,Book>} & Account role K and owning book. Cross-book effects
require separate balanced batches linked by an economic event.\\
\code{Observation<T>} & Source assertion with contextual and precision
metadata; accessing its value retains its constraint and provenance.\\
\code{Symbol<T,D>} & Shared unknown with domain D; repeated references denote
the same value in a history.\\
\code{Choice<E>} & A finite structural choice over enumeration E. No probability
is implied.\\
\code{Distribution<T>} & Explicit probability measure with parameter,
conditioning and measure metadata. It is not a Symbol or a Choice.\\
\code{SourceRef}, \code{RuleRef}, \code{ModelRef} & Versioned evidence or
implementation references; unresolved references prevent concrete execution.\\
\code{Measurement<T>} & Value plus valuation date, basis, method, assumptions
and numerical status. Recognition remains a rule decision.\\
\bottomrule
\end{longtable}

The portable profile spells the rate type \code{Rate<convention>} and makes
both calendar and accrual convention explicit in
\code{YearFraction<calendar, convention>}. \code{InstantValue<T>} and
\code{Flow<T, [start,end)>} distinguish dated stocks from interval flows;
their context carries the calendar. An approximate numerical component of a
historical \code{Measurement<T>} uses \code{Approximation} with its method,
error description and environment references. An exact measurement remains
exact; valuation context and recognition authority are separate in either case.

Core expressions use exact integers and rationals, Boolean logic, enumerations,
dates, records, lists, comparisons and registered pure functions. A decimal
literal is parsed exactly, not first converted through binary floating point.
Numerical plugins can return approximate measurements, but must declare their
error or numerical status and the rounding used when posting money.

Missingness is a tagged type with distinct values: \code{notDisclosed},
\code{notApplicable}, \code{unreadable}, \code{withheld}, and
\code{unresolvedMeaning}. Zero is an ordinary known number. A missing datum can
be replaced by a declared Symbol or assumption in a model branch; the source
observation stays missing. An unknown is not necessarily random, and an
estimate is not necessarily missing.
On portable import, \code{notDisclosed}, \code{unreadable}, \code{withheld} and
\code{unresolvedMeaning} become \code{Unknown} with the corresponding reason
and obligation references. \code{notApplicable} becomes the explicit
\code{NotApplicable} variant, never an unknown or zero.

Addition of monetary observations also requires compatible entity, perimeter,
time, dimensions and basis. These checks operate on contextual expressions in
addition to the currency type. An analyst can deliberately aggregate dimensions
or transform a basis using an explicit function. For example,
\code{sumOver(segment, ...)} records the dimension being removed and requires
the segment partition to be declared; it does not assume every disclosed
segment total is additive.

\section{A portable record has an exact meaning}
\label{sec:portable-record}

A constructor and a textual frontend can describe the same financial model
only if they exchange values whose identity is independent of a Python object
or an incidental JSON writer. The proposed portable record therefore has six
fields: \code{kind}, \code{schema}, \code{key}, \code{revision}, \code{body} and
\code{dependencies}. Kind and schema identify a registered version; key is its
stable, schema-scoped logical name; revision is a nonnegative integer. The body obeys that
schema's required, optional and nullable fields. Unknown fields reject.

A typed reference has exactly kind, schema and digest. The digest covers the
complete target record. Every body reference is also listed in dependencies;
additional dependencies need a schema-declared role. The submitted closure is
finite and acyclic. A receipt can refer to its request and old state, but the
old state cannot refer forward to that receipt. Stable business names can
relate counterparties without creating a cycle of content references.

Financial quantities use a reduced rational with signed numerator text $n$ and
positive denominator text $d$, with $\gcd(|n|,d)=1$ and zero represented by
$0/1$. The decimal token 2.50 therefore denotes $5/2$ exactly. Revisions,
ordinals and integer counts remain integers; a Boolean is never an integer.
Approximate measurements instead retain decimal text, method, error description
and environment. A rounded decimal cannot silently acquire exact status.

An absent required field is invalid. Unknown financial content is an explicit
value with a reason and outstanding evidence references. It can be retained,
but cannot satisfy a method that requires the missing fact. Null is permitted
only where the particular schema gives it a meaning; not applicable, unknown
and zero are distinct. A candidate also differs from a value admitted in a
specific immutable view. Admission checks the actual decision and its ancestors;
wrapping a candidate in an ``admitted'' type is insufficient.

Portable bytes use UTF-8 JSON with ASCII object keys sorted by their byte order,
no insignificant whitespace, no floats, duplicate keys, nonfinite numbers or
invalid Unicode scalar values. Decoded strings retain their code points.
An escaped Unicode character and the same literal character encode equally;
composed and decomposed spellings remain different unless an explicit named
transformation relates them. Global NFC normalisation would silently change
some source text and is therefore not part of this encoding.

The writer emits literal non-ASCII scalars, escapes quotation marks and reverse
solidus, and escapes control characters U+0000--U+001F. It uses the five named
JSON escapes where defined and lowercase \code{\textbackslash u00xx} for the
remaining controls; solidus is not escaped. Integers use minimal base-ten
spelling without plus signs, leading zeroes or negative zero. Dependencies,
module declarations by key, dimensions by axis and explicit all-of/any-of
operand sets are sorted after rejecting duplicates. Other arrays, including
journal lines, event schedules and diagnostic priority, retain order.

With these rules, the record digest is
\[
\operatorname{SHA256}\bigl(
 \operatorname{UTF8}(\text{``accdsl.portable.v1/''}+\text{schema})
 \;\Vert\;\texttt{0x00}\;\Vert\;\operatorname{canonical}(\text{record})\bigr).
\]
This recognises representational identity, not economic equivalence or source
truth. Original documents retain a separate hash of their raw bytes. A blob
reference names that byte hash, algorithm and length; it is not a typed record
reference. A submission record links those bytes to media type, transport and
decoder, preserving a malformed request without pretending it was valid.

Historical Core 0 and replay examples keep their original encodings and IDs.
A declared legacy import retains original bytes, decoder, old identity, portable
output and field-by-field comparison. If one historical record maps to several
portable records, an import mapping names every output and conversion. Import
does not grant financial admission in a new view. This separation permits old
examples to remain reproducible while the proposed portable representation
becomes precise enough for two implementations to exchange.

\section{The first company profile and its authoring boundary}
\label{sec:portable-authoring}

The proposed \code{CorporateJointHistory/1} profile connects one legal entity and
book per branch to annual/interim sources, revisions, working capital, payroll,
tangible assets, simple financing, an explicitly declared tax model and complete
statements. Its types preserve entity, perimeter and currency even where an
operation is outside this initial scope. A real company cannot become eligible
for the profile by dropping its subsidiaries, leases or contingencies.

Typed constructors must express every included declaration in the portable IR
without opaque Python functions. The initially matched text frontend covers the
bounded debt and constraint fragment. Whole-company textual authoring remains
a subsequent frontend obligation; the whole-company typed model is part of the
initial domain. Thus domain scope and authoring syntax are separate decisions.
The following Core 0 listing teaches the shared debt family under its preserved
historical grammar. It is not a promise that the first new parser supports every
historical declaration.
The matched portable debt example uses floating share $x$, while the following
listing uses fixed share $\alpha$. The substitution $x=1-\alpha$ preserves
the entire family, including its endpoint results.

Backend admission examines expanded expressions. Multiplication of an unknown
quantity by an unknown price is bilinear even when hidden inside a function
called revenue. The initial exact reconstruction fragment admits rational affine
constraints, finite choices and bounded integers under explicit templates.
Unsupported functions return a capability diagnostic. A type-correct historical
representation may therefore be retained even when a requested forecast still
requires an unavailable method or an additional admitted assumption.

The portable \code{CompileRequest} supplies source kind, source bytes or a typed
module, schema and profile references, the full package lock, evidence/time view
and requested capability target. Its stages are strict decoding and source
locations; name/import resolution; dependency and cycle checks; static type and
context checks; admission-layer checks; expanded-operation capability analysis;
and canonical record production. The earlier six-stage workflow describes the
historical Core 0 frontend; it does not add execution planning to this portable
compiler or permit it to call a solver.

The result contains candidate IR, source map, unresolved inputs, capability
analysis and a validation receipt bound to every input. Invalid structure or
types yield no executable IR. Valid structure with an unresolved financial input
can yield a retained candidate whose affected target has execution admission
false. Source maps retain decoded Unicode code-point half-open ranges for text
or constructor paths for typed input. Formatting can change source-map identity
without changing financial-model identity.

For a matched example, both authoring routes must describe debt 100, floating
share $x\in[0,1]$ and conditional interest
$100((1-x)/20+7x/100)=5+2x$. Endpoints are 5 and 7; $x=1/2$ gives 6. Changing
the floating assumption from 7\% to 9\% changes the upper endpoint to 9.
Independent arithmetic is needed because two frontends could share the same
wrong constant. Agreement between them alone does not establish correctness.

\section{A small model with a continuous interpretation family}

The following source uses a synthetic observation whose full context is stored
in the evidence registry. The example imports the demonstration debt package,
which uses one-year simple accrual, no fees and immediate rate resets. It is
not an IFRS rule package.

All dependencies are instantiated in
\code{fixtures/demo-debt-package.json}: the source reading, balanced opening
state, currency precision, contract state and dates, reset and interest
operations, closing and report mappings, request and expected outputs.
\code{fixtures/borrowing.acc} is the complete source, including the retained
earnings account and an \code{opening} declaration. The shorter listings below
explain its declarations in stages. They do not replace that full source.

\begin{lstlisting}
module BorrowingExample version "1";
use DemoDebt version "1";
source note7 = external("synthetic-note-7");
entity Acme functional USD;
book Local entity Acme basis DemoDebt version "1";
account Cash : Asset book Local;
account Debt : Liability book Local;
account InterestExpense : Expense book Local;

observation debt : Money<USD> instant @2025-12-31
  = 100 USD from note7;
symbol alpha : Real in [0, 1]
  because "The fixed-rate proportion is not disclosed";
\end{lstlisting}

The remaining declarations use this same unknown fraction in both loan
components. A change to the fraction therefore changes the whole debt mix
consistently.

\begin{lstlisting}
let fixedPrincipal = alpha * value(debt);
let floatingPrincipal = (1 - alpha) * value(debt);
constraint split:
  fixedPrincipal + floatingPrincipal == value(debt);

contract loans : DebtMix book Local {
  principal = value(debt);
  fixed_fraction = alpha;
  fixed_rate = 0.05 per year;
  initial_floating_rate = 0.05 per year;
}
alternative allFixed { assume alpha == 1; }
alternative allFloating { assume alpha == 0; }
scenario rateShock {
  set market_rate = constant(0.07 per year);
}
\end{lstlisting}

The parent model retains every $\alpha\in[0,1]$. The two named alternatives
are useful conditional views of that family, not a declaration that the
endpoints exhaust it. Requesting the parent family's exact bounds produces
$[5,7]$ for the next year's expense under \code{rateShock}. Requesting a
conditional run of \code{allFixed} supplies $\alpha=1$. Requesting a weighted
mixture requires a separate weight model; the two labels do not create equal
weights.

An observation declaration is a reference to evidence plus a proposed typed
reading. The compiler checks that the registry context agrees with the declared
period and type. If the evidence is rounded, \code{value(debt)} denotes a
shared constrained value rather than a silently exact 100. This synthetic
registry explicitly declares an exact value, which keeps the illustration
deterministic.

\begin{minipage}{\linewidth}
\section{Posting rules and explicit exceptions}

An event type declares the information needed by a rule. The example below
assumes an event representing interest accrued and paid on the same date.
Separate accrual and settlement events are needed when those dates differ.

\begin{lstlisting}
event InterestPaid(principal: Money<USD>,
                   rate: Rate<annual_simple>,
                   fraction: YearFraction) book Local;

rule recogniseInterest for InterestPaid as e
when e.principal >= 0 USD {
  let charge = roundMoney(simpleInterest(e.principal, e.rate, e.fraction),
                         currencyScale(USD), "half_even");
  require charge >= 0 USD;
  post {
    debit InterestExpense charge;
    credit Cash charge;
  }
}
\end{lstlisting}
\end{minipage}

The rule cannot perform file I/O, mutate another entity or sample a random
number. It returns a proposed batch. The runtime verifies balance and all
preconditions before committing state. Negative interest requires an explicit
supported rule; it is not mishandled by this intentionally restricted example.
Rounding is a declared demonstration convention, not a universal accounting
default.

Rules can declare \code{overrides ruleName}. The resulting override graph
must be acyclic. For a concrete event, the evaluator determines all applicable
rules and removes any that are overridden by another applicable rule. Exactly
one maximal rule must remain for a required recognition operation. An empty
set of applicable rules produces \code{UnsupportedTreatment}; more than one
maximal rule produces
\code{RuleConflict}. File order and a hidden numeric priority never break ties.

If a guard depends on a symbolic value, the solver partitions the current
feasible family into the guard and its negation, subject to the existing
constraints. This is a logical partition, not a random draw. An incomplete
solver returns \code{UnknownApplicability}, not a guessed rule. Statically
provable disjointness or total coverage can be cached; otherwise the runtime
performs the same checks on every concrete event.

\section{Alternative readings, assumptions and updates}

An alternative adds constraints and possibly a supported semantic mapping to
its parent model. Its record contains a rationale, supporting spans, opposing
evidence, author, creation time and scope. Child alternatives inherit all
parent constraints; they cannot quietly discard an inconvenient disclosure.
A deliberate source rejection is a separate review decision with an explicit
reason and a new evidence view.

Two independent alternatives may coexist, while two mutually exclusive readings
of the same text require an explicit choice group. The representation stores
these relationships in a constraint graph. It does not take a Cartesian
product blindly. For example, interpreting an instrument as a liability affects
both its balance-sheet account and the treatment of its periodic payments;
those choices must be linked through one shared treatment variable.

New evidence narrows an existing family by conjunction when it concerns the
same definitions and reporting basis. If it changes the basis or corrects an
earlier assertion, the system creates a revised evidence view and records the
relationship. Posterior reweighting, when enabled, is a separate operation
with likelihood provenance. A branch that becomes inconsistent is retained
as historically considered and rejected, with the evidence that rejected it.

\section{Modules and stable intermediate representation}

A package identifies its namespace, semantic version, dependencies, effective
dates, jurisdiction, supported entity classes, recognition operations and
measurement policies. Standard packages define allowable treatments; industry
packages add economic objects and specialised rules; company packages map
source concepts and declare evidenced elections. A company extension cannot
override a mandatory standard rule merely because its package is loaded last.
An override is valid only at a declared extension point and within the selected
standard's policy domain.

The IR has separate node classes for evidence, contextual observations,
expressions, domains, constraints, alternatives, contracts, economic events,
measurement results, rules, journal batches, report mappings and runs.
Each node has a stable identifier and explicit input references. No expression
stores executable host-language text. Plugins are referenced by versioned
identifiers with typed signatures. This permits a second implementation to
execute the same supported program without evaluating arbitrary embedded code.

The IR serialisation version is independent of the DSL version and of any
accounting-standard edition. Migrations preserve the original payload and
record changed meanings. A syntax migration can be automatic when semantics
are unchanged. A changed recognition policy requires a new model version,
even if the source file still parses.

\begin{minipage}{\linewidth}
\section{Errors are part of the interface}

The compiler and runtime return typed diagnostics with a source span,
affected entity and period, dependency chain and suggested local action.
The principal diagnostic categories have distinct recovery paths:

\begin{center}
\begin{tabular}{L{61mm}L{75mm}}
\toprule Diagnostic & Meaning and next action\\\midrule
\code{TypeMismatch}, \code{ContextMismatch} & Resolve incompatible units,
periods, entities or measurement bases; do not coerce implicitly.\\
\code{MissingEvidence} & Supply a source or label a model input as an assumption.\\
\code{InfeasibleHistory} & Inspect the conflicting evidence and model
constraints; preserve the rejected candidate.\\
\code{RuleConflict}, \code{UnknownApplicability} & Resolve rule scope or
preserve a justified conditional partition.\\
\code{UnsupportedTreatment} & Add a reviewed capability or leave the affected
output incomplete.\\
\code{UnbalancedBatch} & Reject the atomic update and diagnose the posting rule.\\
\code{SolverIncomplete}, \code{NumericalFailure} & Preserve the target and
inputs; repair or change the computational method.\\
\code{InsufficientSimulationState} & Add declared conditional assumptions or
request a family of results.\\\bottomrule
\end{tabular}
\end{center}

Warnings can describe low extraction confidence or an approximate economic
model, but they do not replace these errors. A required unknown input prevents
a concrete unconditional result. The analyst can add a declared conditional
assumption, request bounds, or leave that output incomplete. Every path has a
different, explicit result type; there is no generic fallback to zero.
\end{minipage}\par
